\chapter{高级算法专题：数据结构、数学、字符串和图论}

前面的章节已经覆盖了面试高频主线：数组、哈希、二分、链表、栈队列、堆、树图、回溯、动态规划、贪心和常见 C++ 容器。本章补齐更高阶的算法版图。它不是为了堆砌竞赛名词，而是回答一个更根本的问题：当普通数组、哈希表、BFS、DFS 和基础 DP 已经不够时，下一层优化从哪里来。

高级算法通常来自五类结构性观察。第一，查询很多而修改较少，可以预处理，例如稀疏表、倍增 LCA、前缀和、Floyd。第二，查询和修改都在线发生，需要能局部更新的结构，例如树状数组和线段树。第三，状态空间太大，但 key 可以拆成字符、二进制位或图分量，例如 Trie、AC 自动机、后缀结构和强连通缩点。第四，暴力转移枚举太多前驱，需要单调队列、分治、斜率或离线排序降低复杂度。第五，问题本身带有数学结构，例如同余、素数、组合数、矩阵乘法、几何向量和概率抽样。

\begin{principlebox}{高级专题的学习顺序}
\begin{tabular}{@{}p{0.23\linewidth}p{0.69\linewidth}@{}}
先问模型 & 这是动态点查询、区间查询、路径查询、字符串匹配、图连通性、流量分配，还是数学计数。\\
再问暴力 & 暴力到底慢在哪里，是重复扫描区间、重复搜索路径、重复匹配模式串，还是重复枚举状态。\\
然后找性质 & 可加性、单调性、可分解性、前缀可维护、边权非负、图可缩点、集合可压缩，都可能是优化来源。\\
最后写代码 & 代码只是性质的落地。实现时必须讲清下标、边界、内存、溢出、迭代器和复杂度。\\
\end{tabular}
\end{principlebox}

本章的定位是“面试和算法教材之间的高级桥梁”。对于普通后端、客户端、基础设施面试，树状数组、线段树、KMP、Dijkstra、并查集、拓扑、背包、状态压缩、快速幂和组合数已经很有价值；对于更难岗位或竞赛型笔试，Tarjan、网络流、后缀结构、数位 DP、单调队列优化、几何和随机结构会继续拉开差距。你不必第一遍背完所有代码，但必须知道每种算法解决什么问题、为什么正确、什么时候不该用。

\section{高级数据结构：区间、顺序和离线查询}

树状数组解决的是一类非常稳定的问题：数组会在线修改，同时需要查询前缀和或可由前缀差得到的区间和。暴力方法每次区间查询扫描 \code{[l,r]}，单次 \complexity{O(n)}；如果查询很多就会慢。前缀和能把查询降到 \complexity{O(1)}，但单点修改会影响后面所有前缀，单次修改又变成 \complexity{O(n)}。树状数组的优化点是把前缀拆成若干个二进制低位决定的块，每次修改只影响包含该位置的少量块，每次查询只累加覆盖前缀的少量块。

\begin{mentalmodel}
树状数组不是一棵显式指针树，而是一个数组上的二进制分块结构。 \code{tree[i]} 管理的区间长度是 \code{lowbit(i)}，也就是 \code{i & -i}。查询前缀时不断去掉最低位一，修改时不断加上最低位一。每一步都改变一个二进制块，所以复杂度是 \complexity{O(\log n)}。
\end{mentalmodel}

\begin{lstlisting}[language=C++,caption={树状数组：单点增加和前缀求和}]
class FenwickTree {
public:
    explicit FenwickTree(int n) : tree_(n + 1, 0) {}

    void add(int index, long long delta)
    {
        for (int i = index + 1; i < static_cast<int>(this->tree_.size());
             i += i & -i) {
            this->tree_[i] += delta;
        }
    }

    long long prefix_sum(int count) const
    {
        long long result = 0;
        for (int i = count; i > 0; i -= i & -i) {
            result += this->tree_[i];
        }
        return result;
    }

    long long range_sum(int left, int right) const
    {
        return this->prefix_sum(right + 1) - this->prefix_sum(left);
    }

private:
    std::vector<long long> tree_;
};
\end{lstlisting}

这个实现对外使用 0 基下标，对内使用 1 基下标。 \code{prefix_sum(count)} 表示前 \code{count} 个元素的和，也就是区间 \code{[0,count)}。这样 \code{range_sum(left,right)} 就能写成前缀差，不容易混淆。树状数组常见案例包括 LeetCode 307 区域和检索、LeetCode 315 计算右侧小于当前元素的个数、LeetCode 493 翻转对。后两题还需要离散化，因为值域可能很大，树状数组管理的是排名而不是原始数值。

以“右侧小于当前元素的个数”为例，暴力做法对每个位置向右扫描，复杂度 \complexity{O(n^2)}。从右往左处理时，已经加入树状数组的都是当前位置右侧元素；要知道有多少数小于 \code{nums[i]}，只需查询排名小于它的元素频次。优化来源不是神秘技巧，而是把“右侧历史元素集合”维护成一个支持排名前缀查询的数据结构。

\begin{lstlisting}[language=C++,caption={LeetCode 315：离散化加树状数组统计右侧较小元素}]
std::vector<int> count_smaller(const std::vector<int>& nums)
{
    std::vector<int> sorted = nums;
    std::sort(sorted.begin(), sorted.end());
    sorted.erase(std::unique(sorted.begin(), sorted.end()), sorted.end());

    FenwickTree bit(static_cast<int>(sorted.size()));
    std::vector<int> answer(nums.size(), 0);

    for (int i = static_cast<int>(nums.size()) - 1; i >= 0; --i) {
        const int rank = static_cast<int>(
            std::lower_bound(sorted.begin(), sorted.end(), nums[i]) -
            sorted.begin());
        answer[static_cast<std::size_t>(i)] =
            static_cast<int>(bit.prefix_sum(rank));
        bit.add(rank, 1);
    }

    return answer;
}
\end{lstlisting}

线段树比树状数组更通用。树状数组最适合可逆的前缀聚合，例如和；线段树可以保存区间最大值、最小值、和、最大子段和、颜色覆盖、区间赋值、区间加法等更多信息。它的基本思想是把区间递归二分，每个节点保存一个区间的聚合结果。单点修改只影响从叶子到根的一条路径；区间查询把目标区间拆成若干个不相交节点。若还要区间修改，就需要懒标记，把“整段都加了多少”暂存在节点上，等访问子节点时再下推。

\begin{lstlisting}[language=C++,caption={线段树：区间加法和区间最大值查询}]
class LazySegmentTree {
public:
    explicit LazySegmentTree(const std::vector<long long>& values)
        : n_(static_cast<int>(values.size())),
          max_(4 * std::max(1, this->n_), 0),
          lazy_(4 * std::max(1, this->n_), 0)
    {
        if (this->n_ > 0) {
            this->build(1, 0, this->n_ - 1, values);
        }
    }

    void add(int left, int right, long long delta)
    {
        this->add(1, 0, this->n_ - 1, left, right, delta);
    }

    long long query_max(int left, int right)
    {
        return this->query_max(1, 0, this->n_ - 1, left, right);
    }

private:
    void build(int node, int left, int right,
               const std::vector<long long>& values)
    {
        if (left == right) {
            this->max_[node] = values[static_cast<std::size_t>(left)];
            return;
        }
        const int mid = left + (right - left) / 2;
        this->build(node * 2, left, mid, values);
        this->build(node * 2 + 1, mid + 1, right, values);
        this->pull(node);
    }

    void apply(int node, long long delta)
    {
        this->max_[node] += delta;
        this->lazy_[node] += delta;
    }

    void push(int node)
    {
        if (this->lazy_[node] == 0) {
            return;
        }
        this->apply(node * 2, this->lazy_[node]);
        this->apply(node * 2 + 1, this->lazy_[node]);
        this->lazy_[node] = 0;
    }

    void pull(int node)
    {
        this->max_[node] = std::max(this->max_[node * 2],
                                    this->max_[node * 2 + 1]);
    }

    void add(int node, int left, int right,
             int query_left, int query_right, long long delta)
    {
        if (query_left <= left && right <= query_right) {
            this->apply(node, delta);
            return;
        }
        this->push(node);
        const int mid = left + (right - left) / 2;
        if (query_left <= mid) {
            this->add(node * 2, left, mid, query_left, query_right, delta);
        }
        if (mid < query_right) {
            this->add(node * 2 + 1, mid + 1, right,
                      query_left, query_right, delta);
        }
        this->pull(node);
    }

    long long query_max(int node, int left, int right,
                        int query_left, int query_right)
    {
        if (query_left <= left && right <= query_right) {
            return this->max_[node];
        }
        this->push(node);
        const int mid = left + (right - left) / 2;
        long long answer = std::numeric_limits<long long>::lowest();
        if (query_left <= mid) {
            answer = std::max(answer,
                              this->query_max(node * 2, left, mid,
                                              query_left, query_right));
        }
        if (mid < query_right) {
            answer = std::max(answer,
                              this->query_max(node * 2 + 1, mid + 1, right,
                                              query_left, query_right));
        }
        return answer;
    }

    int n_;
    std::vector<long long> max_;
    std::vector<long long> lazy_;
};
\end{lstlisting}

懒标记的正确性来自一个不变量：节点保存的聚合值已经包含所有覆盖到该节点的修改；子节点可能还没有真实更新，但 \code{lazy} 记录了以后下推所需的信息。完全覆盖时只更新当前节点，不继续递归；部分覆盖时必须先下推旧标记，再递归子区间，最后用子节点重新计算当前节点。常见错误是忘记下推就查询子区间，导致子节点值落后；或者下推后忘记清空当前懒标记，导致重复加。

稀疏表适合静态数组的区间最值查询。它不能高效处理在线修改，但预处理后每次查询可以 \complexity{O(1)}。核心是 \code{st[k][i]} 保存从 \code{i} 开始、长度 \code{2^k} 的区间最值。查询 \code{[l,r]} 时，取 \code{k = floor(log2(r-l+1))}，用两个长度 \code{2^k} 的区间覆盖目标区间：一个从 \code{l} 开始，一个以 \code{r} 结尾。最小值、最大值这种幂等操作允许重叠覆盖，所以可以这么做；区间和不能用这种重叠方式直接查询。

\begin{lstlisting}[language=C++,caption={稀疏表：静态区间最小值查询}]
class SparseTableMin {
public:
    explicit SparseTableMin(const std::vector<int>& values)
    {
        const int n = static_cast<int>(values.size());
        this->log_.assign(n + 1, 0);
        for (int i = 2; i <= n; ++i) {
            this->log_[i] = this->log_[i / 2] + 1;
        }

        const int levels = this->log_[n] + 1;
        this->table_.assign(levels, std::vector<int>(n));
        this->table_[0] = values;
        for (int level = 1; level < levels; ++level) {
            const int length = 1 << level;
            for (int i = 0; i + length <= n; ++i) {
                this->table_[level][i] =
                    std::min(this->table_[level - 1][i],
                             this->table_[level - 1][i + length / 2]);
            }
        }
    }

    int range_min(int left, int right) const
    {
        const int length = right - left + 1;
        const int level = this->log_[length];
        return std::min(this->table_[level][left],
                        this->table_[level][right - (1 << level) + 1]);
    }

private:
    std::vector<int> log_;
    std::vector<std::vector<int>> table_;
};
\end{lstlisting}

LCA 最近公共祖先是树上查询的基础问题。暴力方法从两个节点不断向上走，最坏可能 \complexity{O(n)}。若查询很多，可以用倍增预处理。 \code{up[k][v]} 表示节点 \code{v} 向上跳 \code{2^k} 步到达的祖先。查询时先把两个节点跳到同一深度，再从大到小尝试同时跳，直到它们的父节点相同。优化来源是把“一步一步向上”拆成二进制跳跃。

动态区间题还经常用扫描线。扫描线不是一个具体容器，而是一种把二维或时间区间事件排序后按顺序处理的思想。会议室数量、天际线、矩形面积并、航班预订、区间加减都可以写成事件：在左端点加入影响，在右端点移除影响。若事件只需要统计当前数量，用差分或堆即可；若事件需要维护覆盖长度，就需要线段树保存坐标压缩后的区间覆盖次数和实际覆盖长度。核心原则是：把连续变化拆成有限个端点事件，只在事件发生处更新状态。

\section{数学算法：数论、组合和矩阵}

数学算法的第一层是整数边界。很多题看似算法正确，却因为溢出、负数取模或除法方向出错。C++ 中 \code{int} 溢出是未定义行为，乘法取模前应使用 \code{long long} 或更宽类型。负数取模的结果可能仍为负，若要得到半开区间 $[0,\mathrm{mod})$ 范围，应按 $((x \bmod \mathrm{mod}) + \mathrm{mod}) \bmod \mathrm{mod}$ 归一化。这些不是细枝末节，而是数学题和工程代码之间的硬接口。

快速幂把 \code{a^n} 的线性连乘变成二进制分解。暴力乘 \code{n} 次是 \complexity{O(n)}；观察指数的二进制表示，每一位代表是否乘上当前底数，底数每轮平方，指数每轮右移。这个思想也用于矩阵快速幂、排列快速幂和函数复合快速幂。

\begin{lstlisting}[language=C++,caption={快速幂：支持取模的大整数乘方}]
long long mod_pow(long long base, long long exponent, long long mod)
{
    base %= mod;
    long long result = 1 % mod;
    while (exponent > 0) {
        if ((exponent & 1LL) != 0) {
            result = result * base % mod;
        }
        base = base * base % mod;
        exponent >>= 1LL;
    }
    return result;
}
\end{lstlisting}

最大公约数来自欧几里得算法：\code{gcd(a,b) = gcd(b,a mod b)}。它的正确性来自同余：同时整除 \code{a} 和 \code{b} 的数，也会整除 \code{a - q b}；反过来也成立。扩展欧几里得进一步求出整数 \code{x,y}，使 \code{a x + b y = gcd(a,b)}。当 \code{a} 和模数 \code{mod} 互质时，\code{x} 就是 \code{a} 在模 \code{mod} 下的逆元。

\begin{lstlisting}[language=C++,caption={扩展欧几里得：求 Bezout 系数}]
long long extended_gcd(long long a, long long b,
                       long long& x, long long& y)
{
    if (b == 0) {
        x = 1;
        y = 0;
        return a;
    }
    long long next_x = 0;
    long long next_y = 0;
    const long long g = extended_gcd(b, a % b, next_x, next_y);
    x = next_y;
    y = next_x - (a / b) * next_y;
    return g;
}
\end{lstlisting}

素数筛解决批量素性问题。逐个数试除到平方根，若要判断很多数，会重复大量工作。埃氏筛从每个素数出发标记它的倍数，复杂度约 \complexity{O(n \log \log n)}。线性筛保证每个合数只被它的最小素因子标记一次，复杂度 \complexity{O(n)}，但代码和证明更细。普通面试通常埃氏筛足够，竞赛或极限数据再考虑线性筛。

\begin{lstlisting}[language=C++,caption={埃氏筛：生成不超过 n 的全部素数}]
std::vector<int> sieve_primes(int n)
{
    std::vector<char> is_prime(n + 1, true);
    if (n >= 0) {
        is_prime[0] = false;
    }
    if (n >= 1) {
        is_prime[1] = false;
    }

    for (long long p = 2; p * p <= n; ++p) {
        if (!is_prime[static_cast<std::size_t>(p)]) {
            continue;
        }
        for (long long x = p * p; x <= n; x += p) {
            is_prime[static_cast<std::size_t>(x)] = false;
        }
    }

    std::vector<int> primes;
    for (int x = 2; x <= n; ++x) {
        if (is_prime[static_cast<std::size_t>(x)]) {
            primes.push_back(x);
        }
    }
    return primes;
}
\end{lstlisting}

组合数经常出现在路径计数、选子集、概率和动态规划中。如果模数是质数，可以预处理阶乘和逆元，用
\[
C(n,k) = \frac{n!}{k!(n-k)!}
\]
转成乘法逆元。预处理后每次查询 \complexity{O(1)}。若模数不是质数，逆元不一定存在，需要用其他方法；若 \code{n} 极大但 \code{k} 较小，可以逐项乘除并约分；若涉及多个模数，可以用中国剩余定理作为进阶。

\begin{lstlisting}[language=C++,caption={质数模数组合数：阶乘和逆元预处理}]
class CombinationMod {
public:
    CombinationMod(int n, long long mod)
        : mod_(mod), fact_(n + 1, 1), inv_fact_(n + 1, 1)
    {
        for (int i = 1; i <= n; ++i) {
            this->fact_[i] = this->fact_[i - 1] * i % this->mod_;
        }
        this->inv_fact_[n] = mod_pow(this->fact_[n], this->mod_ - 2, this->mod_);
        for (int i = n; i >= 1; --i) {
            this->inv_fact_[i - 1] = this->inv_fact_[i] * i % this->mod_;
        }
    }

    long long choose(int n, int k) const
    {
        if (k < 0 || k > n) {
            return 0;
        }
        return this->fact_[n] * this->inv_fact_[k] % this->mod_ *
               this->inv_fact_[n - k] % this->mod_;
    }

private:
    long long mod_;
    std::vector<long long> fact_;
    std::vector<long long> inv_fact_;
};
\end{lstlisting}

矩阵快速幂适合线性递推。斐波那契数列满足
\[
\begin{bmatrix}F_{n+1}\\F_n\end{bmatrix}
=
\begin{bmatrix}1&1\\1&0\end{bmatrix}
\begin{bmatrix}F_n\\F_{n-1}\end{bmatrix}.
\]
因此可以把第 \code{n} 项转成矩阵的 \code{n} 次幂。更一般地，只要状态转移是固定线性组合，就可以把一次转移写成矩阵乘法，再用快速幂把 \complexity{O(n)} 降到 \complexity{O(k^3 \log n)}，其中 \code{k} 是状态维度。

\begin{lstlisting}[language=C++,caption={矩阵快速幂：固定维度线性递推}]
using Matrix = std::vector<std::vector<long long>>;

Matrix multiply(const Matrix& a, const Matrix& b, long long mod)
{
    const int n = static_cast<int>(a.size());
    const int m = static_cast<int>(b[0].size());
    const int inner = static_cast<int>(b.size());
    Matrix result(n, std::vector<long long>(m, 0));
    for (int i = 0; i < n; ++i) {
        for (int k = 0; k < inner; ++k) {
            if (a[i][k] == 0) {
                continue;
            }
            for (int j = 0; j < m; ++j) {
                result[i][j] =
                    (result[i][j] + a[i][k] * b[k][j]) % mod;
            }
        }
    }
    return result;
}

Matrix matrix_pow(Matrix base, long long exponent, long long mod)
{
    const int n = static_cast<int>(base.size());
    Matrix result(n, std::vector<long long>(n, 0));
    for (int i = 0; i < n; ++i) {
        result[i][i] = 1;
    }
    while (exponent > 0) {
        if ((exponent & 1LL) != 0) {
            result = multiply(result, base, mod);
        }
        base = multiply(base, base, mod);
        exponent >>= 1LL;
    }
    return result;
}
\end{lstlisting}

数学题的复盘不能只写公式。要写清楚变量范围、模数性质、是否存在逆元、乘法是否溢出、负数如何取模、是否需要预处理。LeetCode 50 Pow、LeetCode 204 计数质数、LeetCode 372 超级次方、LeetCode 509 斐波那契、LeetCode 790 多米诺和托米诺平铺，都可以用本节思路串起来。真正难的是把题面转换成同余、递推或组合模型。

\section{字符串算法：从匹配到自动机}

字符串算法的第一层成本是子串。C++ 的 \code{substr} 会创建新字符串，若在双循环里反复调用，复杂度可能从看似 \complexity{O(n^2)} 变成更高。高质量字符串题解要先明确：比较的是字符、前缀、哈希值、自动机状态，还是字典树节点。不要把字符串当作免费切片。

KMP 解决单模式串匹配。暴力匹配在文本每个位置尝试模式串，失败后文本起点右移一位，最坏 \complexity{O(nm)}。KMP 的优化来自模式串自身的 border 信息：当匹配到第 \code{j} 个字符失败时，不必把模式串完全归零，而是跳到当前已匹配前缀的最长真前后缀长度。 \code{pi[i]} 表示模式串 \code{[0,i]} 的最长真前后缀长度。

\begin{lstlisting}[language=C++,caption={KMP：前缀函数和模式串查找}]
std::vector<int> prefix_function(const std::string& pattern)
{
    std::vector<int> pi(pattern.size(), 0);
    for (int i = 1; i < static_cast<int>(pattern.size()); ++i) {
        int j = pi[static_cast<std::size_t>(i - 1)];
        while (j > 0 && pattern[static_cast<std::size_t>(i)] !=
                            pattern[static_cast<std::size_t>(j)]) {
            j = pi[static_cast<std::size_t>(j - 1)];
        }
        if (pattern[static_cast<std::size_t>(i)] ==
            pattern[static_cast<std::size_t>(j)]) {
            ++j;
        }
        pi[static_cast<std::size_t>(i)] = j;
    }
    return pi;
}

std::vector<int> kmp_search(const std::string& text,
                            const std::string& pattern)
{
    std::vector<int> answer;
    if (pattern.empty()) {
        return answer;
    }
    const std::vector<int> pi = prefix_function(pattern);
    int matched = 0;
    for (int i = 0; i < static_cast<int>(text.size()); ++i) {
        while (matched > 0 && text[static_cast<std::size_t>(i)] !=
                                pattern[static_cast<std::size_t>(matched)]) {
            matched = pi[static_cast<std::size_t>(matched - 1)];
        }
        if (text[static_cast<std::size_t>(i)] ==
            pattern[static_cast<std::size_t>(matched)]) {
            ++matched;
        }
        if (matched == static_cast<int>(pattern.size())) {
            answer.push_back(i - matched + 1);
            matched = pi[static_cast<std::size_t>(matched - 1)];
        }
    }
    return answer;
}
\end{lstlisting}

KMP 的正确性来自一个不变量：\code{matched} 始终是当前文本后缀与模式串前缀匹配的最大长度。失配时，任何比 \code{pi[matched-1]} 更长的候选前缀都不可能同时作为已匹配部分的后缀，所以可以直接跳过。LeetCode 28 找出字符串中第一个匹配项、LeetCode 459 重复的子字符串都可以用 KMP 解释。重复子字符串的判断是：若最长 border 长度为 \code{l}，且 \code{n} 能被 \code{n-l} 整除，则字符串由长度 \code{n-l} 的模式重复构成。

Z 函数从另一个角度描述字符串。 \code{z[i]} 表示从 \code{i} 开始的后缀与整个字符串前缀的最长公共前缀长度。它适合处理“每个位置和前缀匹配多少”这类问题。KMP 的 \code{pi} 关注每个前缀的 border，Z 函数关注每个后缀与全局前缀的匹配长度。两者能互相转换，但学习时应分别理解它们的状态含义。

\begin{lstlisting}[language=C++,caption={Z 函数：维护当前最右匹配盒}]
std::vector<int> z_function(const std::string& s)
{
    const int n = static_cast<int>(s.size());
    std::vector<int> z(n, 0);
    int left = 0;
    int right = 0;
    for (int i = 1; i < n; ++i) {
        if (i <= right) {
            z[i] = std::min(right - i + 1, z[i - left]);
        }
        while (i + z[i] < n &&
               s[static_cast<std::size_t>(z[i])] ==
               s[static_cast<std::size_t>(i + z[i])]) {
            ++z[i];
        }
        if (i + z[i] - 1 > right) {
            left = i;
            right = i + z[i] - 1;
        }
    }
    return z;
}
\end{lstlisting}

Manacher 算法求所有中心的最长回文半径。暴力以每个中心向两侧扩展，最坏 \complexity{O(n^2)}。Manacher 的优化和 Z 函数类似：维护当前最右回文区间，利用对称中心的已知半径给当前位置一个初始半径，再继续向外扩展。为了统一奇偶长度，常把原字符串插入分隔符，或分别处理奇回文和偶回文。

\begin{lstlisting}[language=C++,caption={Manacher：奇数长度回文半径}]
std::vector<int> odd_palindrome_radii(const std::string& s)
{
    const int n = static_cast<int>(s.size());
    std::vector<int> radius(n, 0);
    int left = 0;
    int right = -1;
    for (int i = 0; i < n; ++i) {
        int k = 1;
        if (i <= right) {
            const int mirror = left + right - i;
            k = std::min(radius[mirror], right - i + 1);
        }
        while (0 <= i - k && i + k < n &&
               s[static_cast<std::size_t>(i - k)] ==
               s[static_cast<std::size_t>(i + k)]) {
            ++k;
        }
        radius[i] = k;
        if (i + k - 1 > right) {
            left = i - k + 1;
            right = i + k - 1;
        }
    }
    return radius;
}
\end{lstlisting}

Trie 字典树把字符串按字符分层。哈希表适合整串等值查询，Trie 适合前缀查询、自动补全、单词搜索和按字符转移。LeetCode 208 实现 Trie、LeetCode 211 添加与搜索单词、LeetCode 212 单词搜索 II 都属于这个方向。Trie 的代价是节点数可能接近所有单词字符总数；若字符集固定为小写字母，可以用长度 26 的数组；若字符集大或稀疏，用哈希表保存孩子更节省空间但常数更高。

\begin{lstlisting}[language=C++,caption={Trie：小写字母前缀树}]
class Trie {
public:
    void insert(const std::string& word)
    {
        int node = 0;
        for (const char ch : word) {
            const int index = ch - 'a';
            if (this->next_[node][index] == -1) {
                this->next_[node][index] =
                    static_cast<int>(this->next_.size());
                this->next_.push_back({});
                this->next_.back().fill(-1);
                this->terminal_.push_back(false);
            }
            node = this->next_[node][index];
        }
        this->terminal_[node] = true;
    }

    bool starts_with(const std::string& prefix) const
    {
        return this->walk(prefix) != -1;
    }

    bool contains(const std::string& word) const
    {
        const int node = this->walk(word);
        return node != -1 && this->terminal_[node];
    }

    Trie()
    {
        this->next_.push_back({});
        this->next_.back().fill(-1);
        this->terminal_.push_back(false);
    }

private:
    int walk(const std::string& text) const
    {
        int node = 0;
        for (const char ch : text) {
            const int index = ch - 'a';
            if (index < 0 || index >= 26 || this->next_[node][index] == -1) {
                return -1;
            }
            node = this->next_[node][index];
        }
        return node;
    }

    std::vector<std::array<int, 26>> next_;
    std::vector<char> terminal_;
};
\end{lstlisting}

AC 自动机是 Trie 加失败指针，用来多模式串匹配。暴力做法对每个模式串分别 KMP，若模式串很多会重复扫描文本；AC 自动机把所有模式串建成一棵 Trie，再为每个节点建立失败指针，表示当前路径失配时应该跳到哪个最长后缀状态。扫描文本时只走一次自动机，就能同时发现多个模式串。它的本质是把“多个模式串的前缀集合”压成有限状态机。

后缀数组和后缀自动机属于更高级的字符串索引。后缀数组把所有后缀排序，配合 LCP 可以解决最长重复子串、字符串比较、第 k 小子串等问题；后缀自动机压缩了所有子串形成的状态集合，适合统计不同子串数量、最长公共子串和多串匹配。这两者代码较长，普通面试不一定要求手写，但算法教材必须讲清用途：当问题反复询问“所有后缀”“所有子串”“字典序”和“重复结构”时，单次 KMP 或 Trie 已经不够，需要后缀结构。

\section{高级图论：分量、匹配和流}

Tarjan 强连通分量解决有向图中“哪些节点彼此可达”的问题。普通 DFS 可以判断从一个点能到哪里，但无法直接把所有互相可达的节点分组。暴力做法对每个点跑一次 DFS，再两两判断可达性，复杂度很高。Tarjan 在一次 DFS 中维护时间戳 \code{dfn} 和低链接值 \code{low}。 \code{low[u]} 表示从 \code{u} 的 DFS 子树经过零条或多条树边和最多一条返祖边能到达的最早时间戳。当 \code{low[u] == dfn[u]} 时，\code{u} 是一个强连通分量的根，栈中从顶部弹到 \code{u} 的节点形成一个 SCC。

\begin{lstlisting}[language=C++,caption={Tarjan 强连通分量：一次 DFS 缩点}]
class TarjanScc {
public:
    explicit TarjanScc(const std::vector<std::vector<int>>& graph)
        : graph_(graph),
          dfn_(graph.size(), 0),
          low_(graph.size(), 0),
          in_stack_(graph.size(), false)
    {
    }

    std::vector<std::vector<int>> run()
    {
        for (int node = 0; node < static_cast<int>(this->graph_.size());
             ++node) {
            if (this->dfn_[node] == 0) {
                this->dfs(node);
            }
        }
        return this->components_;
    }

private:
    void dfs(int node)
    {
        this->dfn_[node] = this->low_[node] = ++this->timer_;
        this->stack_.push_back(node);
        this->in_stack_[node] = true;

        for (const int next : this->graph_[node]) {
            if (this->dfn_[next] == 0) {
                this->dfs(next);
                this->low_[node] = std::min(this->low_[node],
                                            this->low_[next]);
            } else if (this->in_stack_[next]) {
                this->low_[node] = std::min(this->low_[node],
                                            this->dfn_[next]);
            }
        }

        if (this->low_[node] != this->dfn_[node]) {
            return;
        }
        std::vector<int> component;
        while (true) {
            const int x = this->stack_.back();
            this->stack_.pop_back();
            this->in_stack_[x] = false;
            component.push_back(x);
            if (x == node) {
                break;
            }
        }
        this->components_.push_back(std::move(component));
    }

    const std::vector<std::vector<int>>& graph_;
    std::vector<int> dfn_;
    std::vector<int> low_;
    std::vector<int> stack_;
    std::vector<char> in_stack_;
    std::vector<std::vector<int>> components_;
    int timer_ = 0;
};
\end{lstlisting}

强连通缩点后，原图会变成有向无环图。很多题的关键不是在原图上硬做 DP，而是先把环压成分量，再在 DAG 上拓扑处理。2-SAT、依赖关系、可达性压缩都属于这个思想。要注意：有向图的强连通分量和无向图的连通分量不是一回事；无向图里双向可达天然成立，有向图里必须考虑边方向。

割点和桥处理无向图的脆弱连接。桥是一条删除后会增加连通分量数量的边；割点是删除后会增加连通分量数量的点。它们也使用 DFS 时间戳和低链接值。对于树边 \code{u-v}，若 \code{low[v] > dfn[u]}，说明 \code{v} 的子树无法通过返祖边回到 \code{u} 或更早祖先，边 \code{u-v} 是桥。若 \code{low[v] >= dfn[u]}，则 \code{u} 可能是割点，但根节点需要单独判断子树数量。

二分图匹配解决“左侧对象和右侧对象如何一对一配对最多”的问题。暴力枚举所有匹配组合是指数级。匈牙利算法的核心是增广路：若一个左侧点想匹配某个右侧点，而右侧点已经被别人占用，就尝试让原来的左侧点改配另一个右侧点。只要找到一条从未匹配左点到未匹配右点的交替路径，匹配数就能增加一。

\begin{lstlisting}[language=C++,caption={二分图最大匹配：DFS 增广路}]
int maximum_bipartite_matching(const std::vector<std::vector<int>>& graph,
                               int right_size)
{
    std::vector<int> matched_left(right_size, -1);

    auto try_match = [&](auto&& self, int left,
                         std::vector<char>& seen) -> bool {
        for (const int right : graph[left]) {
            if (seen[right]) {
                continue;
            }
            seen[right] = true;
            if (matched_left[right] == -1 ||
                self(self, matched_left[right], seen)) {
                matched_left[right] = left;
                return true;
            }
        }
        return false;
    };

    int answer = 0;
    for (int left = 0; left < static_cast<int>(graph.size()); ++left) {
        std::vector<char> seen(right_size, false);
        if (try_match(try_match, left, seen)) {
            ++answer;
        }
    }
    return answer;
}
\end{lstlisting}

这里的 \code{seen} 每次从新的左点开始都要清空，它表示本轮增广中右侧点是否尝试过，不是全局访问标记。LeetCode 中类似题包括 LeetCode 785 判断二分图、LeetCode 886 可能的二分法，以及一些任务分配变体。若边权存在且要求最大权匹配，就需要 KM 算法或最小费用流，普通匈牙利只处理无权最大匹配。

网络流把“容量限制下能从源点送到汇点多少流量”形式化。它适合处理多源多汇、选择约束、割、匹配和资源分配。Dinic 算法先用 BFS 建分层图，只允许沿层数增加的边发送流；再用 DFS 在分层图中找阻塞流。残量网络是关键：每条正向边发送流后，反向边增加可撤销容量，表示以后可以取消部分选择，把流改走别的路径。

\begin{lstlisting}[language=C++,caption={Dinic 最大流：分层图和当前弧优化}]
class Dinic {
public:
    struct Edge {
        int to;
        int rev;
        long long cap;
    };

    explicit Dinic(int n) : graph_(n), level_(n), iter_(n) {}

    void add_edge(int from, int to, long long cap)
    {
        Edge forward{to, static_cast<int>(this->graph_[to].size()), cap};
        Edge backward{from, static_cast<int>(this->graph_[from].size()), 0};
        this->graph_[from].push_back(forward);
        this->graph_[to].push_back(backward);
    }

    long long max_flow(int source, int sink)
    {
        long long flow = 0;
        while (this->bfs(source, sink)) {
            std::fill(this->iter_.begin(), this->iter_.end(), 0);
            while (true) {
                const long long pushed =
                    this->dfs(source, sink, std::numeric_limits<long long>::max());
                if (pushed == 0) {
                    break;
                }
                flow += pushed;
            }
        }
        return flow;
    }

private:
    bool bfs(int source, int sink)
    {
        std::fill(this->level_.begin(), this->level_.end(), -1);
        std::queue<int> queue;
        this->level_[source] = 0;
        queue.push(source);
        while (!queue.empty()) {
            const int node = queue.front();
            queue.pop();
            for (const Edge& edge : this->graph_[node]) {
                if (edge.cap > 0 && this->level_[edge.to] < 0) {
                    this->level_[edge.to] = this->level_[node] + 1;
                    queue.push(edge.to);
                }
            }
        }
        return this->level_[sink] >= 0;
    }

    long long dfs(int node, int sink, long long flow)
    {
        if (node == sink) {
            return flow;
        }
        for (int& i = this->iter_[node];
             i < static_cast<int>(this->graph_[node].size()); ++i) {
            Edge& edge = this->graph_[node][i];
            if (edge.cap <= 0 ||
                this->level_[node] + 1 != this->level_[edge.to]) {
                continue;
            }
            const long long pushed =
                this->dfs(edge.to, sink, std::min(flow, edge.cap));
            if (pushed == 0) {
                continue;
            }
            edge.cap -= pushed;
            this->graph_[edge.to][edge.rev].cap += pushed;
            return pushed;
        }
        return 0;
    }

    std::vector<std::vector<Edge>> graph_;
    std::vector<int> level_;
    std::vector<int> iter_;
};
\end{lstlisting}

网络流代码比普通图搜索长，但思想仍然来自暴力改进。暴力尝试一条可行路径送流，可能反复走很差路径；残量网络允许撤销，分层图限制每一阶段只走更接近汇点的边，当前弧避免重复扫描已经无效的边。二分图最大匹配可以转成最大流：源点连左侧点，左侧点连右侧点，右侧点连汇点，所有容量为一。这样匹配、割和容量限制进入同一个框架。

\section{高级动态规划、搜索和离线算法}

数位 DP 处理“统计不超过 N 的数中有多少满足某种数字性质”。暴力枚举所有数再检查，若 N 很大不可行。数位 DP 把上界按十进制或二进制拆成位，从高位到低位构造数字。状态通常包括当前位置、是否贴着上界、是否已经开始、以及与题目性质相关的统计量。例如统计不含连续一的非负整数，可以用二进制位、上一位是否为一、是否受上界限制来定义状态。

\begin{lstlisting}[language=C++,caption={LeetCode 600：数位 DP 统计不含连续一的整数}]
int find_integers_without_consecutive_ones(int n)
{
    std::vector<int> bits;
    for (int x = n; x > 0; x >>= 1) {
        bits.push_back(x & 1);
    }
    std::reverse(bits.begin(), bits.end());

    std::map<std::tuple<int, int, int>, int> memo;
    auto dfs = [&](auto&& self, int pos, int previous_one,
                   int tight) -> int {
        if (pos == static_cast<int>(bits.size())) {
            return 1;
        }
        const auto key = std::make_tuple(pos, previous_one, tight);
        if (!tight && memo.contains(key)) {
            return memo[key];
        }
        const int limit = tight ? bits[static_cast<std::size_t>(pos)] : 1;
        int answer = 0;
        for (int bit = 0; bit <= limit; ++bit) {
            if (previous_one && bit == 1) {
                continue;
            }
            answer += self(self, pos + 1, bit, tight && bit == limit);
        }
        if (!tight) {
            memo[key] = answer;
        }
        return answer;
    };

    return dfs(dfs, 0, 0, 1);
}
\end{lstlisting}

数位 DP 的 \code{tight} 表示当前前缀是否和上界完全相同。只有 \code{tight == false} 的状态才能安全记忆化，因为它的后续选择不再依赖具体上界前缀。若题目要求正整数，还要处理前导零；若题目要求数字和、模数、出现次数，就把这些量加入状态。高质量讲解必须说明每个状态维度的必要性，不能把模板参数机械列出。

折半搜索处理的是 $2^n$ 太大但 $2^{n/2}$ 可接受的场景。子集和、最接近目标的子序列和、四数组求和都可用这个思想。把数组分成两半，分别枚举所有子集和，再通过排序、二分或双指针合并。暴力枚举完整集合是 \complexity{O(2^n)}；折半后是 \complexity{O(2^{n/2} \log 2^{n/2})} 量级。优化来源不是剪枝，而是把全集选择分成两个独立半区。

\begin{lstlisting}[language=C++,caption={LeetCode 1755：折半搜索最接近目标的子序列和}]
std::vector<int> subset_sums(const std::vector<int>& nums,
                             int left, int right)
{
    std::vector<int> sums{0};
    for (int i = left; i < right; ++i) {
        const int current_size = static_cast<int>(sums.size());
        for (int j = 0; j < current_size; ++j) {
            sums.push_back(sums[static_cast<std::size_t>(j)] + nums[i]);
        }
    }
    return sums;
}

int min_abs_difference(const std::vector<int>& nums, int goal)
{
    const int n = static_cast<int>(nums.size());
    std::vector<int> left_sums = subset_sums(nums, 0, n / 2);
    std::vector<int> right_sums = subset_sums(nums, n / 2, n);
    std::sort(right_sums.begin(), right_sums.end());

    int answer = std::abs(goal);
    for (const int left_sum : left_sums) {
        const int need = goal - left_sum;
        auto it = std::lower_bound(right_sums.begin(), right_sums.end(), need);
        if (it != right_sums.end()) {
            answer = std::min(answer, std::abs(left_sum + *it - goal));
        }
        if (it != right_sums.begin()) {
            --it;
            answer = std::min(answer, std::abs(left_sum + *it - goal));
        }
    }
    return answer;
}
\end{lstlisting}

IDA* 和 A* 都是带启发式的搜索。普通 BFS 保证最短步数，但状态爆炸；DFS 省空间但不保证最短。A* 用 \code{f = g + h} 排序，\code{g} 是已走代价，\code{h} 是对剩余代价的低估。若 \code{h} 从不高估真实剩余代价，A* 第一次弹出目标时能保证最优。IDA* 则用深度优先搜索配合逐渐增加的估价上界，常见于十五数码这类状态空间大但启发式较强的问题。面试中不一定手写它们，但要知道启发式必须可采纳，不能为了跑得快写一个会高估的 \code{h} 后还声称最优。

离线算法把所有查询先收集起来，重新排序后统一处理。它适合“答案不要求按输入过程实时输出”的问题。区间中小于等于某值的元素个数，可以把数组元素按值排序、查询按阈值排序，用树状数组逐步加入可用元素；并查集也常用于离线连通性，把边按权重或时间排序后逐步合并。离线的本质是改变处理顺序，用单调推进替代每个查询从头扫描。

莫队算法是离线区间查询的代表。若每个查询 \code{[l,r]} 的答案可以在移动左右端点时 \complexity{O(1)} 或接近 \complexity{O(1)} 更新，就可以把查询按块排序，让相邻查询的端点移动距离变小。它常用于静态数组上的区间不同数、区间频次、区间众数近似处理等。莫队不是万能：如果加入和删除无法快速维护，或者存在修改操作但没有扩展成带修莫队，就不适合。

\section{几何、随机结构和工程算法}

计算几何的基本单位是向量。点积回答夹角和投影，叉积回答方向和面积。给三个点 \code{a,b,c}，向量 \code{b-a} 与 \code{c-a} 的叉积为正表示 \code{c} 在有向线段 \code{a->b} 的左侧，为负表示右侧，为零表示共线。线段相交、凸包、多边形面积都建立在这个判断上。整数坐标下，叉积可能溢出，应使用 \code{long long}。

\begin{lstlisting}[language=C++,caption={几何基础：叉积和线段相交}]
struct Point {
    long long x;
    long long y;
};

long long cross(Point a, Point b, Point c)
{
    const long long x1 = b.x - a.x;
    const long long y1 = b.y - a.y;
    const long long x2 = c.x - a.x;
    const long long y2 = c.y - a.y;
    return x1 * y2 - y1 * x2;
}

bool between(long long a, long long b, long long x)
{
    return std::min(a, b) <= x && x <= std::max(a, b);
}

bool segments_intersect(Point a, Point b, Point c, Point d)
{
    const long long c1 = cross(a, b, c);
    const long long c2 = cross(a, b, d);
    const long long c3 = cross(c, d, a);
    const long long c4 = cross(c, d, b);

    if (c1 == 0 && between(a.x, b.x, c.x) && between(a.y, b.y, c.y)) {
        return true;
    }
    if (c2 == 0 && between(a.x, b.x, d.x) && between(a.y, b.y, d.y)) {
        return true;
    }
    if (c3 == 0 && between(c.x, d.x, a.x) && between(c.y, d.y, a.y)) {
        return true;
    }
    if (c4 == 0 && between(c.x, d.x, b.x) && between(c.y, d.y, b.y)) {
        return true;
    }
    return (c1 > 0) != (c2 > 0) && (c3 > 0) != (c4 > 0);
}
\end{lstlisting}

凸包求一组点的最外层边界。Andrew 单调链算法先按坐标排序，再分别构造下凸壳和上凸壳。每加入一个新点，若最后两个点和新点构成非左转，就弹出中间点，因为它不可能在凸包边界上。这个过程和单调栈非常像：维护一个满足几何凸性的候选栈。LeetCode 587 安装栅栏就是凸包题，边界点是否保留要根据题意调整弹出条件。

随机算法首先要保证分布正确。蓄水池抽样解决“数据流长度未知，从中均匀随机选一个或 k 个元素”。选一个元素时，第 \code{i} 个元素以 \code{1/i} 的概率替换当前答案。证明用归纳：第 \code{i} 轮后，任意历史元素保留概率都是 \code{1/i}。这个算法只用 \complexity{O(1)} 空间，适合流式数据。

\begin{lstlisting}[language=C++,caption={蓄水池抽样：数据流中均匀选择一个元素}]
template <class T, class URBG>
T reservoir_sample_one(std::istream& input, URBG& rng)
{
    T answer{};
    T value{};
    int count = 0;
    while (input >> value) {
        ++count;
        std::uniform_int_distribution<int> dist(1, count);
        if (dist(rng) == 1) {
            answer = value;
        }
    }
    return answer;
}
\end{lstlisting}

布隆过滤器用多个哈希函数和一个位数组判断“某个元素是否可能出现过”。插入时把多个哈希位置置一；查询时若有任一位置为零，则元素一定没出现；若全部为一，则元素可能出现。它没有假阴性，但有假阳性。适合缓存穿透防护、去重预判和大规模集合近似查询。它不能删除元素，除非使用计数布隆过滤器；它也不能告诉你元素出现次数。

一致性哈希用于分布式系统中的节点扩缩容。普通取模 $hash(key) \bmod n$ 在节点数量变化时会让大量 key 重新映射；一致性哈希把节点和 key 都放到环上，key 归属顺时针遇到的第一个节点。增加或删除节点时，只影响环上相邻区间。虚拟节点用于缓解负载不均。这个结构不属于传统力扣题高频内容，但它展示了哈希思想在工程系统中的另一种使用方式：目标不是单机 \complexity{O(1)} 查询，而是降低重分布成本。

跳表用多层随机索引实现有序集合。底层是完整有序链表，上层以概率抽样提升部分节点，查询时从最高层向右走，走不动就下沉。期望复杂度 \complexity{O(\log n)}。与红黑树相比，跳表实现更直接，天然适合范围扫描；Redis 的有序集合底层就使用跳表和哈希表组合。它的正确性不是依赖严格平衡，而是依赖随机层高带来的期望高度。

这些工程算法提醒我们：算法教材不能只看最坏复杂度，还要看数据分布、内存布局、并发、更新频率和错误概率。布隆过滤器允许假阳性换空间，一致性哈希允许少量重映射换扩缩容稳定，跳表允许随机化换实现简洁，蓄水池抽样允许单遍扫描换均匀样本。面试表达中能说清这些取舍，比机械背代码更接近工程能力。

\section{高级专题深讲：从题面到算法选择}

高级题最怕两种错误。第一种是看到关键词就套模板，例如看到“区间”就写线段树，看到“字符串”就写 KMP，看到“图”就写 Dijkstra。第二种是只会最终代码，不知道为什么能从暴力降下来。真正的算法选择应该从题面的操作集合出发：数据是否静态，查询是否很多，修改是否在线，答案是否要求最优路径，状态是否可以重复到达，比较对象是单点、区间、集合、前缀、后缀还是路径。

\begin{principlebox}{算法选择的四问}
\begin{tabular}{@{}p{0.2\linewidth}p{0.72\linewidth}@{}}
对象是什么 & 数组元素、区间、字符串前缀、图节点、边、集合状态、数字位、几何点。\\
操作是什么 & 查询、修改、插入、删除、合并、匹配、计数、求最优、判断可达。\\
暴力慢在哪 & 重复扫描、重复匹配、重复搜索、重复枚举前驱、重复比较路径。\\
性质是什么 & 可加、可逆、幂等、单调、可排序、可缩点、可分块、可二进制拆分。\\
\end{tabular}
\end{principlebox}

区间类题可以先画一条时间轴。若题目只给一次数组，之后问很多静态区间最值，稀疏表比线段树更简单；若需要单点修改和区间和，树状数组通常比线段树更轻；若需要区间加法和区间最大值，线段树的懒标记才必要；若所有操作都能离线排序，扫描线或树状数组离线处理可能比在线结构更清楚。结构越强，代码越重，不要用重结构掩盖简单性质。

以 LeetCode 731 我的日程安排 II 为例，题目允许预订半开区间，要求任意时刻最多双重预订，不能出现三重预订。暴力做法是把每次预订后所有时间点都扫描一遍，这在时间值很大时不可行。第一种结构化做法是差分扫描：每次临时把开始点加一、结束点减一，按时间顺序累计，若累计超过二就回滚。这说明本题本质是“动态事件点上的前缀最大值”。若预订次数不大，\code{std::map} 差分已经足够；若次数巨大且要极致性能，可以考虑动态线段树维护区间最大覆盖次数。

\begin{lstlisting}[language=C++,caption={LeetCode 731：有序差分检查三重预订}]
class MyCalendarTwo {
public:
    bool book(int start, int end)
    {
        ++this->diff_[start];
        --this->diff_[end];

        int active = 0;
        for (const auto& [time, delta] : this->diff_) {
            active += delta;
            if (active > 2) {
                this->rollback(start, end);
                return false;
            }
        }
        return true;
    }

private:
    void rollback(int start, int end)
    {
        if (--this->diff_[start] == 0) {
            this->diff_.erase(start);
        }
        if (++this->diff_[end] == 0) {
            this->diff_.erase(end);
        }
    }

    std::map<int, int> diff_;
};
\end{lstlisting}

这段代码的正确性依赖半开区间语义。一个预订在 \code{start} 时刻生效，在 \code{end} 时刻失效，所以两个区间首尾相接不重叠。扫描差分时，\code{active} 表示当前时间段内的预订数量；如果它超过二，说明刚刚加入的区间制造了三重覆盖，必须撤销。这个版本单次操作是 \complexity{O(m)}，其中 \code{m} 是不同端点数。它不一定是最优复杂度，但推导清晰，非常适合面试第一版。若面试官继续追问性能，再说明动态线段树可以把每次操作降到 \complexity{O(\log U)}，其中 \code{U} 是时间值域。

LeetCode 699 掉落的方块体现了另一个区间问题：每个方块落到一个横坐标区间上，新高度等于该区间已有最大高度加方块边长，然后整个区间被赋成新高度。暴力做法对每个方块扫描之前所有方块判断重叠，复杂度 \complexity{O(n^2)}。若 \code{n} 较小，这个暴力加区间判断可以接受；若要系统优化，就把所有端点离散化，用线段树维护区间最大值和区间赋值。这里的关键是“坐标值很大但端点数量有限”，所以离散化不是可选细节，而是把巨大值域压成可维护数组的前提。

LCA 的选型也要从查询次数出发。若只问一次二叉树最近公共祖先，递归返回目标节点很直接；若一棵静态树上有大量任意节点 LCA 查询，倍增预处理更合适；若还要支持边权路径最大值或最小值，可以在倍增表旁边再维护跳跃路径上的聚合值。倍增的本质不是“背表”，而是把向上走很多步拆成二进制跳跃。每个查询先把深节点跳到同一深度，再从高位到低位同时尝试跳过不相同的祖先。

\begin{lstlisting}[language=C++,caption={树上倍增：迭代建表和 LCA 查询}]
class BinaryLiftingLca {
public:
    explicit BinaryLiftingLca(const std::vector<std::vector<int>>& tree,
                              int root)
        : depth_(tree.size(), 0)
    {
        int levels = 1;
        while ((1 << levels) <= static_cast<int>(tree.size())) {
            ++levels;
        }
        this->up_.assign(levels, std::vector<int>(tree.size(), root));

        std::queue<int> queue;
        std::vector<char> visited(tree.size(), false);
        queue.push(root);
        visited[root] = true;
        this->up_[0][root] = root;

        while (!queue.empty()) {
            const int node = queue.front();
            queue.pop();
            for (const int next : tree[node]) {
                if (visited[next]) {
                    continue;
                }
                visited[next] = true;
                this->depth_[next] = this->depth_[node] + 1;
                this->up_[0][next] = node;
                queue.push(next);
            }
        }

        for (int level = 1; level < levels; ++level) {
            for (int node = 0; node < static_cast<int>(tree.size()); ++node) {
                this->up_[level][node] =
                    this->up_[level - 1][this->up_[level - 1][node]];
            }
        }
    }

    int lca(int a, int b) const
    {
        if (this->depth_[a] < this->depth_[b]) {
            std::swap(a, b);
        }
        int diff = this->depth_[a] - this->depth_[b];
        for (int level = 0; diff > 0; ++level) {
            if ((diff & 1) != 0) {
                a = this->up_[level][a];
            }
            diff >>= 1;
        }
        if (a == b) {
            return a;
        }
        for (int level = static_cast<int>(this->up_.size()) - 1;
             level >= 0; --level) {
            if (this->up_[level][a] != this->up_[level][b]) {
                a = this->up_[level][a];
                b = this->up_[level][b];
            }
        }
        return this->up_[0][a];
    }

private:
    std::vector<int> depth_;
    std::vector<std::vector<int>> up_;
};
\end{lstlisting}

这段代码用 BFS 建深度和父节点，避免深树递归爆栈。查询时先处理深度差，再从最高层向下跳。若两个节点在某一层的祖先不同，说明 LCA 在这些祖先之上，可以同时跳过去；若相同，则不能跳，否则会越过答案。最后两个节点停在 LCA 的不同子节点上，返回它们的父亲。LeetCode 1483 树节点的第 K 个祖先正是倍增表的直接应用；LeetCode 236 若只有一次查询，不必上倍增，但理解倍增能帮助你处理多查询版本。

\section{字符串专题深讲：自动机、哈希和后缀结构}

字符串算法的分水岭是“匹配一次”还是“建立索引后反复查询”。KMP 和 Z 函数适合单个模式串或围绕一个字符串的前缀关系；Trie 适合许多单词的前缀查询；AC 自动机适合许多模式串同时在一篇文本里出现；滚动哈希适合快速比较子串；后缀数组和后缀自动机适合全局后缀、全局子串和重复结构。把这些算法混在一起背，会导致选型混乱。

滚动哈希的暴力来源很直观：比较两个长度为 \code{len} 的子串，逐字符需要 \complexity{O(len)}。若要反复比较大量子串，就可以把前缀哈希预处理好，使任意子串哈希值 \complexity{O(1)} 取出。常用多项式哈希：
\[
H[i+1] = H[i] \cdot base + s[i].
\]
子串哈希由前缀差得到。哈希有碰撞风险，工程中可以用双模数、64 位自然溢出或最终回退到字符串比较。教材里必须强调：哈希相等只表示大概率相等，不是数学绝对证明。

\begin{lstlisting}[language=C++,caption={滚动哈希：常数时间取子串哈希}]
class RollingHash {
public:
    explicit RollingHash(const std::string& text)
        : hash_(text.size() + 1, 0), power_(text.size() + 1, 1)
    {
        constexpr unsigned long long BASE = 1315423911ULL;
        for (int i = 0; i < static_cast<int>(text.size()); ++i) {
            this->hash_[i + 1] = this->hash_[i] * BASE +
                                 static_cast<unsigned char>(text[i]) + 1ULL;
            this->power_[i + 1] = this->power_[i] * BASE;
        }
    }

    unsigned long long range_hash(int left, int right) const
    {
        return this->hash_[right] -
               this->hash_[left] * this->power_[right - left];
    }

private:
    std::vector<unsigned long long> hash_;
    std::vector<unsigned long long> power_;
};
\end{lstlisting}

LeetCode 1044 最长重复子串可以用“二分长度 + 滚动哈希”推导。暴力枚举长度和所有子串再比较，成本很高。若固定长度 \code{len}，问题变成是否存在两个相同子串，可以把所有长度为 \code{len} 的子串哈希放入集合。若存在长度为 \code{len} 的重复子串，那么更短长度也必然存在重复，所以长度具有单调性，可以二分。这里优化由两层性质组成：哈希降低子串比较成本，单调性降低长度枚举成本。

AC 自动机可以从 Trie 的失败处理推导出来。假设模式串集合是 \code{he}、\code{she}、\code{his}、\code{hers}。文本扫描到某个字符失配时，不能简单回到根，因为当前已经匹配的后缀可能正好是另一个模式串的前缀。失败指针指向“当前路径字符串的最长可用真后缀状态”。构建时用 BFS 从浅到深处理，因为一个节点的失败指针依赖父节点的失败转移。

\begin{mentalmodel}
AC 自动机的状态不是某个单词，而是“当前文本后缀等于 Trie 中哪条前缀”。普通 Trie 只知道沿成功边往下走；失败指针告诉你失配时保留哪个最长后缀。输出链或累计输出数量告诉你当前状态对应哪些模式串已经出现。
\end{mentalmodel}

LeetCode 212 单词搜索 II 是 Trie 与回溯的结合，不需要 AC 自动机。题目是在网格上走路径，每次只能向相邻格子移动，文本不是一条固定线性串，所以多模式线性匹配的 AC 自动机并不贴合。优化点是 Trie 前缀剪枝：回溯路径形成的前缀若不在 Trie 中，后面再走也不可能形成任何单词，可以立即停止。这个案例说明高级算法不能只看“很多单词”，还要看文本结构：线性文本用 AC 自动机，网格路径用 Trie 剪枝。

后缀数组可以用“把每个后缀当成字符串排序”来理解。若要找最长重复子串，排序后相似的后缀会相邻，答案就在相邻后缀的最长公共前缀中。暴力排序后缀直接比较字符串会很慢；倍增算法用长度为 $2^k$ 的排名来推出长度为 $2^{k+1}$ 的排名。后缀自动机则从另一侧建模：它把所有子串压进一台自动机，每个状态表示一组 endpos 相同的子串。后缀自动机更抽象，但在统计不同子串数量时非常简洁：每个状态新增的不同子串数量是 \code{len[state] - len[link[state]]}。

\begin{principlebox}{字符串算法选型表}
\begin{tabular}{@{}p{0.25\linewidth}p{0.67\linewidth}@{}}
单模式匹配 & KMP、Z 函数、滚动哈希。\\
很多模式匹配线性文本 & Trie 加 AC 自动机。\\
网格里找很多单词 & Trie 加 DFS 回溯剪枝，不是 AC 自动机。\\
大量子串相等比较 & 滚动哈希，必要时双哈希或回退比较。\\
所有后缀、字典序、最长重复 & 后缀数组和 LCP。\\
所有子串计数、最长公共子串 & 后缀自动机。\\
\end{tabular}
\end{principlebox}

\section{图论专题深讲：建模比算法名字更重要}

图论难题往往输在建模，而不是输在不会写队列或堆。题目给的对象可能是单词、密码、网格、账户、课程、字符串状态、棋盘局面、城市、服务器或任务；只有把“状态是什么，边是什么，边权是什么，目标是什么”说清楚，算法才有落脚点。普通 BFS、Dijkstra、并查集、Tarjan、匹配和网络流的差别，都是在回答不同的图问题。

LeetCode 778 水位上升的泳池中游泳有三种典型建模。第一种是 Dijkstra：进入一个格子的代价是路径上最大高度，扩展邻居时代价取 \code{max(current, height[next])}。第二种是二分答案：猜一个时间 \code{t}，只允许走高度不超过 \code{t} 的格子，检查是否可达。第三种是并查集：按高度从低到高开放格子，每开放一个格子就与已开放邻居合并，起点和终点第一次连通时的高度就是答案。这三种写法都正确，但证明不同。Dijkstra 依赖路径代价单调不降；二分依赖可达性随时间单调；并查集依赖从小到大开放时第一次连通的最小性。

\begin{principlebox}{同一道图题的三种证明}
\begin{tabular}{@{}p{0.2\linewidth}p{0.72\linewidth}@{}}
Dijkstra & 每次确定当前最小体力状态，未来绕路不会降低已经弹出的最大边权。\\
二分答案 & 若时间 t 可达，则任意更大的时间都可达，所以可行性单调。\\
并查集 & 按高度开放节点，第一次让起终点连通的高度就是所有路径最大高度的最小值。\\
\end{tabular}
\end{principlebox}

LeetCode 1192 查找集群内的关键连接是桥的直接应用。暴力方法删除每条边后跑一次 DFS，检查图是否仍连通，复杂度 \complexity{O(m(n+m))}。Tarjan 桥算法一次 DFS 就能完成，因为它记录每个子树是否存在返祖边。如果从 \code{v} 子树无法回到 \code{u} 或 \code{u} 的祖先，那么边 \code{u-v} 就是唯一连接这棵子树和外界的通道。实现时最好给无向边编号，跳过“进入当前节点的那条边”，而不是只跳过父节点；否则重边会被误判。

差分约束是最短路的另一种建模。若约束形如 \code{x_j - x_i <= c}，可以建边 \code{i -> j} 权重 \code{c}，表示 \code{x_j <= x_i + c}。一组约束是否可行，可以转成图中是否存在负环；求最大或最小可行值，可以通过最短路或最长路处理。很多排程、相对顺序、上下界问题都能这样转化。它提醒我们：图的边不一定来自题面里的道路，也可以来自不等式。

网络流建模时要抓住“容量守恒”。如果题目要求从若干资源中选一些分配给若干需求，每个资源有容量，每个需求有上限，中间有允许关系，就可以考虑源点、资源层、需求层、汇点。二分图匹配是容量都为一的特殊流；最小割常对应“切断哪些关系使代价最小”；费用流则在容量之外给每单位流增加成本。普通面试不一定要求手写费用流，但看到“容量 + 成本 + 最优分配”时要知道它属于流模型。

\section{动态规划专题深讲：状态、转移和优化边界}

高级 DP 的核心不是状态数量多，而是状态语义更抽象。基础 DP 的状态通常是位置或容量；高级 DP 会把集合、区间、树节点、数字位、概率、博弈局面和优化结构放进状态。写 DP 前先写一句话定义状态，定义里必须包含“处理到哪里”和“保留什么信息”。如果这句话含糊，代码再短也不可靠。

树形 DP 的状态通常表示“当前子树给父节点提供什么信息”。LeetCode 337 打家劫舍 III 是最清楚的例子。暴力递归对每个节点选择偷或不偷；若偷当前节点，孩子不能偷；若不偷当前节点，孩子可偷可不偷。重复子问题来自同一子树被祖父和父节点的不同选择反复计算。状态可以定义为一对值：\code{take} 表示偷当前节点时当前子树最大收益，\code{skip} 表示不偷当前节点时当前子树最大收益。

\begin{lstlisting}[language=C++,caption={LeetCode 337：树形 DP 的两个返回值}]
struct RobState {
    int take = 0;
    int skip = 0;
};

RobState rob_tree_iterative(TreeNode* root)
{
    if (root == nullptr) {
        return {};
    }

    std::vector<std::pair<TreeNode*, bool>> stack{{root, false}};
    std::unordered_map<TreeNode*, RobState> dp;
    while (!stack.empty()) {
        const auto [node, visited] = stack.back();
        stack.pop_back();
        if (node == nullptr) {
            continue;
        }
        if (!visited) {
            stack.emplace_back(node, true);
            stack.emplace_back(node->right, false);
            stack.emplace_back(node->left, false);
            continue;
        }
        const RobState left = dp[node->left];
        const RobState right = dp[node->right];
        RobState current;
        current.take = node->val + left.skip + right.skip;
        current.skip = std::max(left.take, left.skip) +
                       std::max(right.take, right.skip);
        dp[node] = current;
    }
    return dp[root];
}
\end{lstlisting}

这里用显式栈模拟后序遍历，避免深树递归。 \code{take} 和 \code{skip} 的语义必须互斥且完整：父节点只需要知道子树在“根偷”和“根不偷”两种约束下的最优值，不需要知道子树内部具体偷了哪些节点。很多树形 DP 都遵循这个模式：先问当前节点给父节点留下什么约束，再决定返回几个值。

概率 DP 和期望 DP 要先定义随机过程。普通最短路求最小步数，概率题求成功概率或期望次数。若状态之间有概率转移，期望满足线性方程：当前期望等于一步成本加后继期望的加权平均。简单有向无环状态可以按拓扑顺序 DP；有环状态可能需要方程组或利用吸收马尔可夫链思想。面试里常见的是骰子概率、骑士概率、汤问题这类有限状态 DP。关键不是公式复杂，而是每个转移概率是否加起来为一，边界状态是否吸收。

博弈 DP 则要明确双方都最优。常用状态是“当前局面先手能否必胜”或“当前玩家相对另一个玩家的最大分差”。LeetCode 486 预测赢家可以定义 \code{dp[left][right]} 为当前玩家在区间内能比对手多拿多少分。若选左端，剩下区间轮到对手，所以净优势是 \code{nums[left] - dp[left+1][right]}；选右端同理。最后 \code{dp[0][n-1] >= 0} 表示先手不输。这个定义比直接模拟两个人分数更稳，因为它把“轮到谁”压进了相对分差。

斜率优化和分治优化属于“不要乱用”的高级技巧。它们要求转移式具有特定数学形态。斜率优化通常来自
\[
dp[i] = \min_j \{dp[j] + a_i b_j + c_j\}
\]
这样的直线查询；分治优化要求最优决策点单调；四边形不等式优化要求代价函数满足更强结构。若题目只是普通 \code{min} 转移，没有证明这些性质，套优化会得到错误答案。教材中必须把这些技巧放在“条件满足时使用”的位置，而不是把它们包装成万能模板。

\section{高级算法实验：把抽象结构跑出来}

高级算法需要实验，不只需要读题解。树状数组实验应打印每个下标管理的区间。对一个长度为 16 的数组，列出 \code{tree[1]} 到 \code{tree[16]} 分别覆盖哪些原数组位置，再手动执行一次 \code{add(6, delta)}，观察哪些块被更新。这样 \code{i += i & -i} 就不再是咒语，而是沿着包含该点的更大块向上走。

线段树实验应同时打印节点区间、节点值和懒标记。先对区间做一次加法，再只查询其中一半，观察懒标记什么时候下推。很多错误只有在“先整段修改，再局部查询，再局部修改，再整段查询”的操作序列中才会暴露。测试样例不要只查刚刚修改过的完整区间，要故意查跨越多个节点的区间。

字符串实验应比较四种方法的成本。给一个长文本和很多模式串，分别测试逐模式 KMP、Trie 网格剪枝、AC 自动机线性扫描和滚动哈希比较子串。记录每种方法的适用输入：KMP 适合一个模式，AC 适合多个模式和线性文本，Trie 适合前缀和网格路径，哈希适合大量子串相等比较。实验目标不是跑分，而是让选型和输入形状绑定起来。

图论实验应把同一道题用不同模型做一遍。以最小体力路径为例，分别实现 Dijkstra、二分 BFS、并查集开放节点。对同一个网格打印每个算法的中间状态：Dijkstra 的堆顶代价，二分的可行性结果，并查集的开放高度。你会看到三种算法在证明上完全不同，但都围绕同一个目标：最小化路径上的最大边权。

DP 实验应强制写暴力搜索和记忆化版本。先写指数级递归，打印状态参数；再加 memo，观察重复状态数量；最后改成表格或滚动数组。这样能避免“直接背表格”的问题。背包、区间 DP、数位 DP、状态压缩都适合这套流程。每次优化都必须回答：状态是否少了，转移是否快了，还是只是写法换了。

\begin{labbox}
高级专题实验报告格式：题目和编号；暴力状态或枚举对象；暴力复杂度；可利用性质；所选结构或算法；关键不变量；C++ 边界；三组最小反例；最终复杂度；如果输入条件变化，算法是否失效。每一题按这个格式写，比单纯刷十题更有效。
\end{labbox}

\section{案例清单和复盘路径}

本章补充的算法很多，复习时不要按名字背。建议按问题类型建立索引。动态区间查询先问是否有修改：无修改可考虑前缀和、稀疏表、排序离线；有单点修改和前缀聚合考虑树状数组；有区间修改、区间最大最小或复杂聚合考虑线段树。树上路径查询先问是否静态：静态多查询考虑 LCA、倍增、欧拉序和 RMQ；动态树再进入更高级结构，普通面试很少要求。

字符串题先问是单模式、多模式、前缀还是所有子串。单模式匹配用 KMP 或 Z 函数；多模式匹配用 Trie 加 AC 自动机；回文半径用 Manacher；所有后缀和所有子串的排序、计数、重复结构才考虑后缀数组或后缀自动机。不要拿 KMP 解所有字符串题，也不要一看到字符串就上哈希。哈希适合快速比较，自动机适合状态转移，后缀结构适合全局子串集合。

图论题先问边权和目标。只问连通性，优先并查集或 DFS；无权最短路用 BFS；边权非负用 Dijkstra；有负边或限制边数用 Bellman-Ford；所有点对且节点少用 Floyd；互相可达分组用强连通；无向图脆弱边点用桥和割点；一对一配对用二分图匹配；容量分配和割用网络流。每换一个模型，都要说明普通 BFS/DFS 为什么不够。

动态规划题先问状态是否包含位置、容量、集合、区间、树节点、数字位或概率。背包看物品能用几次和方案是否区分顺序；区间 DP 看最后一步是切分还是选择最后处理对象；状态压缩看 \code{n} 是否允许 \code{2^n}；数位 DP 看是否受上界限制和前导零；优化 DP 看转移是否有窗口最值、决策单调、凸性或可离线。没有证明条件时，不要强行套优化。

\begin{exercisebox}
按下面顺序做一轮高级专题训练：树状数组做 307、315、493；线段树做 731、732、699；稀疏表和 LCA 做 236、1483；数学做 50、204、372、790；字符串做 28、459、214、208、212；图论做 1192、210、743、778、785、886；匹配和流至少用小样例手写建图；DP 做 312、516、600、698、1425、1755。每道题都写三段复盘：暴力慢在哪里，优化保存了什么信息，边界为什么不漏。
\end{exercisebox}

高级算法的最终目标不是把书变成代码模板库，而是让读者面对陌生题时能主动归类。题目要求区间反复修改查询，就想到分块、树状数组、线段树或离线；题目要求多个模式同时匹配，就想到 Trie 和自动机；题目要求互相可达和依赖压缩，就想到强连通分量；题目要求从巨大上界中统计数字，就想到数位 DP；题目要求容量和割，就想到网络流。算法越高级，越要先讲清它解决的单点问题。否则名字越多，理解越散。
